\section{Transition reconstruction and label reconciliation}
\label{app:transitions}

\paragraph{Label orientation.} The estimation output orders regimes by the raw
fitted index, in which state $0$ carries the higher diffusion coefficient and
occupies roughly $67\%$ of the sample with a mean duration of $7.56$ hours,
while state $1$ carries the lower coefficient and a mean duration of $3.76$
hours. Our convention is the opposite, with index $0$ denoting the calm regime.
Adopting it requires swapping the volatilities and, because the logistic
coefficients are indexed by the same labels, swapping the transition
coefficients crosswise: the production $\gamma_{01}$, governing calm-to-stress
transitions, is the raw $\gamma_{10}$, and conversely. Failing to apply the
swap to both objects would produce a model in which volatilities and transition
dynamics refer to different states, which is a silent error rather than one that
any single diagnostic would reveal. We enforce the ordering with an invariant
check requiring $\sigma^y_0 < \sigma^y_1$ at load time.

The orientation is worth a remark on its own terms. A high-volatility state that
is both dominant and longer-lived inverts the usual reading in which stress is
rare and brief. Whether this reflects genuine market structure at hourly
frequency, or the labelling of a transform-space variance decomposition that
does not correspond to economic stress, we cannot determine from the available
artefacts, and we use the terms calm and stress throughout as labels for
variance states rather than as economic claims.

\paragraph{Intercept reconstruction.} Slope coefficients and average marginal
effects are supplied by the estimation output, but intercepts are not. The mean
departure rate from state $i$ over the covariate distribution must match the
reported mean duration $d_i$, which gives
\begin{equation}
  \E_z\!\left[\Lambda\!\left(\alpha_{01}+\gamma_{01}z\right)\right]
     = \frac{1}{d_{\text{calm}}},
  \qquad
  \E_z\!\left[\Lambda\!\left(\alpha_{10}+\gamma_{10}z\right)\right]
     = \frac{1}{d_{\text{stress}}},
\end{equation}
for $\alpha_{01}$ and $\alpha_{10}$ separately, taking the expectation over the
full standardised covariate sample and using a monotone bracketing root finder.
The target rates are $0.2660$ and $0.1323$ per hour. The resulting pair
reproduces the reported stationary stress share of $0.6746$ to within $1.0\%$
under an ergodic approximation and $4.5\%$ under simulation along the realised
covariate path.

This is moment matching, not estimation. The reconstructed intercepts have no
sampling distribution, no standard errors and no likelihood interpretation, and
any inference conditioned on their precise values would be unfounded. They are
reported as Derived in Table~\ref{tab:provenance} for that reason.
